%~~~~~~~~~~~~~~~~~~~~~~~~~~~~~%
%  dernière                   %
%  modification = 09/04/2023  %
%~~~~~~~~~~~~~~~~~~~~~~~~~~~~~%
%      Colas Bardavid


% ===================================================================== %
% ============================ Size pages ============================= %
% \def\setSizeOfPages{%
% \geometry{	inner				= 2.25cm,
% 			outer				= 2.5cm,
% 			top					= 1.75cm,
% 			bottom				= 2.75cm,
% 			footskip			= 9mm,
% 			%showframe
% 		 }
% }

% \setSizeOfPages

% ===================================================================== %
% ======================== Force displaystyle ========================= %
\everymath{\displaystyle}


% ===================================================================== %
% ============================ left/right ============================= %
% see: https://tex.stackexchange.com/a/2610/8323
\let\originalleft\left
\let\originalright\right
\renewcommand{\left}{\mathopen{}\mathclose\bgroup\originalleft}
\renewcommand{\right}{\aftergroup\egroup\originalright}


% ===================================================================== %
% ============================= Spacing =============================== %
\setlength{\parindent}{0pt}
\setlength{\parskip}{0.15cm}


% ===================================================================== %
% ========================= The colored box =========================== %
\newtcolorbox{coloredBox}		   {colback 		= \couleurFond,
									colframe 		= \couleurFond,
									boxsep 			= 0pt,
									top 			= 2mm,
									bottom 			= 2mm,
									left			= 3mm,
									right			= 3mm,
									arc				= 1.5mm,
									width			= 0.7\linewidth,
									center}

% ===================================================================== %
% ========================= The colored box =========================== %
\newtcolorbox{coloredBoxLarge}		{colback 		= \couleurFond,
									colframe 		= \couleurFond,
									boxsep 			= 0pt,
									top 			= 2mm,
									bottom 			= 2mm,
									left			= 3mm,
									right			= 3mm,
									arc				= 1.5mm,
									width			= 0.9\linewidth,
									center}

% ===================================================================== %
% ====================== Box for shuffled answers ===================== %
\newtcolorbox{boxForShuffledAnswers}{colback 		= white,
									colframe 		= \couleurFoncee,
									boxsep 			= 0pt,
									top 			= 2mm,
									bottom 			= 2mm,
									left			= 3mm,
									right			= 3mm,
									arc				= 1.5mm,
									width			= \linewidth,
									boxrule			= 1mm,
									halign			= center,
									title			= Réponses mélangées,
									halign title	= center,
									fonttitle		= \sffamily\bfseries,
									before upper    = {\setstretch{1.9}},
									center}


% ===================================================================== %
% ==================== Box for title corrige fiche ==================== %
\newtcolorbox{boxForTitleCorrigeFiche}{colback 		= \couleurFond,
									colframe 		= \couleurFond,
									boxsep 			= 0pt,
									top 			= 2mm,
									bottom 			= 2mm,
									left			= 10mm,
									right			= 3mm,
									arc				= 0.5mm,
									width			= \linewidth,
									boxrule			= 1mm,
									center}


% ===================================================================== %
% ============================ Tight dotfill ========================== %
% https://tex.stackexchange.com/a/85338/8323
\makeatletter
\newcommand \tightdotfill {\leavevmode \cleaders \hb@xt@ .33em{\hss .\hss }\hfill \kern \z@}
\makeatother


% ===================================================================== %
% =================== Space settings for columns ====================== %
\setlength{\columnsep}{7mm}
% see https://tex.stackexchange.com/a/37919/8323
\setlength{\multicolsep}{5.0pt plus 2.0pt minus 1.5pt}% 50% of original values



% ===================================================================== %
% =============================== Footer ============================== %
\pagestyle{fancy}

\def\myHeaderAndFooterStyle{
\renewcommand{\headrulewidth}{0pt}%
\renewcommand{\footrulewidth}{0.4pt}%
\fancyhf{}
\fancyhf[OLF]{{\footnotesize \textForFooter}}
\fancyhf[ORF]{{\footnotesize \thepage}}
\fancyhf[ERF]{{\footnotesize \textForFooter}}
\fancyhf[ELF]{{\footnotesize \thepage}}}

% \fancypagestyle{plain}{\myHeaderAndFooterStyle}
% \myHeaderAndFooterStyle

\def\startNoFooter{\global\def\textForFooter{}}
\def\startFooterForSolutions{\global\def\textForFooter{Réponses et corrigés}}
\def\startFooterForFiches{\global\def\textForFooter{Fiche \no\printCounterForSheet{}. \varForTitreFicheEntrainement}}
\startFooterForFiches


% ===================================================================== %
% ======================= Table des matières ========================== %

% \contentsmargin{0cm}

% \titlecontents*{section}[0em]
%   {}
%   {\\[0mm]\makebox[0pt][l]{\thecontentslabel\textbf{.}}\hspace{12mm}}
%   {}
%   {\dotfill  \thecontentspage\\[2mm]}%\\[-3mm]}
%   [][]

% \setcounter{tocdepth}{1}     % Dans la table des matieres

% \addto\captionsfrench{% Replace "english" with the language you use
%   \renewcommand{\contentsname}%
%     {Sommaire}%
% }


% ===================================================================== %
% ========= Gestion des parties dans la table des matières ============ %

% \newcommand{\ajoutePartieTableDesMatieres}[1]
% {\addtocontents{toc}{\hrule\prevdepth=0pt\par\vspace*{1.5mm}{\strut\large \bfseries #1}\smallskip\par}}


% ===================================================================== %
% ================== Nice symbol for \leq and \geq ==================== %

\renewcommand\leq{\leqslant}
\renewcommand\geq{\geqslant}
\renewcommand\le{\leqslant}
\renewcommand\ge{\geqslant}


% ===================================================================== %
% ========================== Managing QCM ============================= %

% see: https://tex.stackexchange.com/a/354024/8323

\RequirePackage{tikz}

\newcommand\QCMcircled[1]{\tikz[baseline=(char.base)]{%
        \node[circle,draw,minimum size=1.25em,inner sep=0] (char) {#1};}}
\newcommand\QCMcircledalt[1]{\vphantom{\LARGE A}\tikz[baseline=(char.base)]{%
        \node[circle,draw,inner sep=0pt] (char) {\strut #1};}}

\newcommand{\circleForAnswer}[1]{\protect\QCMcircled{{#1}}}

\newlist{listeQCMaux}{enumerate}{1}
\setlist[listeQCMaux]{	label 		= \protect\QCMcircledalt{{\alph*}},
						topsep		= 0mm,
						partopsep	= 0mm,
						nolistsep}

\newenvironment{listeQCM1Colonne}{\begin{listeQCMaux}}{\end{listeQCMaux}\vspace{-5mm}}
\newenvironment{listeQCM2Colonnes}{\begin{multicols}{2}\begin{listeQCMaux}}{\end{listeQCMaux}\end{multicols}\vspace{-5mm}}
\newenvironment{listeQCM3Colonnes}{\begin{multicols}{3}\begin{listeQCMaux}}{\end{listeQCMaux}\end{multicols}\vspace{-5mm}}
\newenvironment{listeQCM4Colonnes}{\begin{multicols}{4}\begin{listeQCMaux}}{\end{listeQCMaux}\end{multicols}\vspace{-5mm}}
\newenvironment{listeQCM5Colonnes}{\begin{multicols}{5}\begin{listeQCMaux}}{\end{listeQCMaux}\end{multicols}\vspace{-5mm}}

\def\reponseA{\circleForAnswer{a}}
\def\reponseB{\circleForAnswer{b}}
\def\reponseC{\circleForAnswer{c}}
\def\reponseD{\circleForAnswer{d}}
\def\reponseE{\circleForAnswer{e}}
\def\reponseF{\circleForAnswer{f}}
\def\reponseX{\circleForAnswer{x}}
\def\reponseY{\circleForAnswer{y}}

% ===================================================================== %
% ================ Managing Image siding some text ==================== %

\newlength{\currentparskip}
% see: https://tex.stackexchange.com/a/43005/8323

\def\varForPourcentageDeLaPartieAGauche{0.4}
\newcommand{\pourcentageDeLaPartieAGauche}[1]{\global\def\varForPourcentageDeLaPartieAGauche{#1}}

\def\initialisationPartieGauche
{\setlength{\currentparskip}{\parskip}%
\begin{minipage}[c][1\totalheight][t]{\varForPourcentageDeLaPartieAGauche\linewidth}%
\setlength{\parskip}{\currentparskip}}

\def\initialisationPartieDroite
{\end{minipage}\hfill
\FPeval{\varForPourcentageDeLaPartieADroite}{clip(1-\varForPourcentageDeLaPartieAGauche)}
\begin{minipage}[c][1\totalheight][t]{\varForPourcentageDeLaPartieADroite\linewidth}%
\setlength{\parskip}{\currentparskip}}

\def\finalisationDuPartageDePage
{\end{minipage}}


% ===================================================================== %
% ============================== Lists ================================ %

% \setlist[itemize, 1]{$\textbullet$, leftmargin = 2em, labelsep = 0.5em}
% \setlist[itemize, 2]{$\rhd$}


% ===================================================================== %
% ======================== Negative spaces ============================ %

\def\minusSmallskip{\vspace*{-\smallskipamount}}
\def\minusMedskip{\vspace*{-\medskipamount}}
\def\minusBigskip{\vspace*{-\bigskipamount}}
